\section{Μοντέλο Μετανάστευσης σε Δακτύλιο}
\label{sec:arch_migration}

Οι ανεξάρτητοι κόμβοι Orchestrator σχηματίζουν έναν ομότιμο (P2P) δακτύλιο μετανάστευσης. Κάθε κόμβος λαμβάνει μια μοναδική 64-bit διεύθυνση δακτυλίου μέσω SHA-256 κατακερματισμού του δικτυακού του URI.

\subsection*{Τυπολογία Δακτυλίου και Ανακάλυψη Διαδόχου}

Η τοπολογία του δακτυλίου ακολουθεί το πρότυπο δακτυλίου Chord \cite{stoica2001chord}: δοθέντος ενός συνόλου $M$ νήσων $\{I_1, I_2, \ldots, I_M\}$, κάθε νήσος $I_j$ αντιστοιχεί σε μια διεύθυνση $h_j = \text{SHA-256}(\text{URI}_j) \bmod 2^{64}$ στον κυκλικό χώρο διευθύνσεων $[0, 2^{64})$. Οι νήσοι ταξινομούνται κατά αύξουσα σειρά διεύθυνσης και ο \emph{διάδοχος} (successor) κάθε νήσου $I_j$ ορίζεται ως η νήσος $I_{(j+1) \bmod M}$ στον ταξινομημένο δακτύλιο. Κατά την εκκίνηση του συστήματος, κάθε Orchestrator εκτελεί ένα βήμα ανακάλυψης (discovery handshake) κατά το οποίο ανακοινώνει το URI του σε όλους τους άλλους γνωστούς κόμβους, ενημερώνοντας τις τοπικές finger tables. Η διαδικασία αυτή διασφαλίζει ότι κάθε Orchestrator γνωρίζει πάντοτε τον άμεσο διάδοχό του στον δακτύλιο, ακόμα και σε περίπτωση δυναμικής προσθήκης νέων νήσων κατά τη διάρκεια της εκτέλεσης.

\subsection*{Πρωτόκολλο Μετανάστευσης}

Ανά τακτό διάστημα γενεών $T_{\text{mig}}$, κάθε νήσος επιλέγει τα κορυφαία $K_{\text{mig}}$ άτομα του πληθυσμού της, δηλαδή τα άτομα με τις υψηλότερες τιμές καταλληλότητας $f(\bm{x})$, και τα σειριοποιεί σε μήνυμα Protocol Buffer. Η σειριοποίηση περιλαμβάνει τόσο τα διανύσματα γονιδίων $\bm{x}$ όσο και τις αντίστοιχες τιμές $f(\bm{x})$, ώστε ο κόμβος παραλαβής να μπορεί να εντάξει τους μετανάστες άμεσα χωρίς επαναξιολόγηση. Η αποστολή εκτελείται \emph{ασύγχρονα} μέσω gRPC unary κλήσης \texttt{MigrateElites} προς τον διάδοχο κόμβο, χωρίς να αποκλείεται η τοπική εξέλιξη. Αυτό εξασφαλίζει ότι η μετανάστευση δεν εισάγει παρατηρήσιμη καθυστέρηση στον κύκλο γενεών.

\subsection*{Ενσωμάτωση Μεταναστών}

Ο κόμβος παραλαβής ενσωματώνει τους $K_{\text{mig}}$ μετανάστες αντικαθιστώντας τα $K_{\text{mig}}$ λιγότερο κατάλληλα άτομα του τοπικού πληθυσμού, στρατηγική γνωστή ως \emph{replacement of worst}. Πιο συγκεκριμένα, εάν $\mathcal{P}$ είναι ο τοπικός πληθυσμός και $\mathcal{E}$ το σύνολο των εισερχόμενων ελίτ μεταναστών, η ενημερωμένη γενεά $\mathcal{P}'$ ορίζεται ως:
\begin{equation}
\label{eq:migration_replace}
\mathcal{P}' = \bigl(\mathcal{P} \setminus \text{BottomK}(\mathcal{P}, K_{\text{mig}})\bigr) \cup \mathcal{E},
\end{equation}
όπου $\text{BottomK}(\mathcal{P}, K_{\text{mig}})$ είναι τα $K_{\text{mig}}$ άτομα με τη χαμηλότερη καταλληλότητα. Η στρατηγική αυτή διατηρεί σταθερό το μέγεθος πληθυσμού $|\mathcal{P}| = \mu$ και εισάγει νέο γενετικό υλικό από γειτονικές νήσους, αποτρέποντας την πρώιμη σύγκλιση σε τοπικά ακρότατα.

\subsection*{Θεωρητική Ανάλυση Σύγκλισης}
 
Η ιδιότητα αντί-πρώιμης σύγκλισης του μοντέλου νήσων τεκμηριώνεται θεωρητικά: η περιοδική μετανάστευση ελίτ μεταξύ νήσων εισάγει γενετική ποικιλομορφία που αναστέλλει τη σύγκλιση σε υποβέλτιστα τοπικά ακρότατα \cite{cantupaz2000parallel}. Ειδικότερα, για δακτύλιο $M$ νήσων, η θεωρία διάδοσης φημών (rumor spreading) των Doerr, Friedrich et al.\ αποδεικνύει ότι ένα ελίτ άτομο που εισάγεται σε μια τυχαία νήσο διαδίδεται σε \emph{όλο} τον δακτύλιο σε $O(\log M)$ γενεές μετανάστευσης \cite{friedrich2017island}. Για $M = 4$ νήσους, αυτό αντιστοιχεί σε μόλις $2{-}3$ γενεές μετανάστευσης, εξασφαλίζοντας ταχύτατη διάχυση ευνοϊκών αλληλομόρφων σε ολόκληρη τη μεταπληθυσμιακή δομή.

\subsection*{Ευαισθησία Παραμέτρων}

Οι κύριοι παράμετροι του μοντέλου μετανάστευσης είναι ο χρόνος μεταξύ μεταναστεύσεων $T_{\text{mig}}$ (σε γενεές) και το μέγεθος μετανάστευσης $K_{\text{mig}}$ (αριθμός ελίτ ατόμων). Μικρές τιμές $T_{\text{mig}}$ προκαλούν υπερβολική ανταλλαγή πληροφοριών και ουσιαστική συγχώνευση των νήσων σε έναν ενιαίο πανμικτικό πληθυσμό, ενώ μεγάλες τιμές μειώνουν την επικοινωνία και δυνητικά αφήνουν νήσους να συγκλίνουν ανεξάρτητα σε διαφορετικά τοπικά ακρότατα. Τα αποτελέσματα της ανάλυσης ευαισθησίας παρουσιάζονται εκτενώς στο Κεφάλαιο~5 και καταλήγουν στις βέλτιστες τιμές $T_{\text{mig}} = 10$ γενεές και $K_{\text{mig}} = 5$ άτομα για τον παρόντα χώρο σχεδίασης.
